\subsection{Chaos-Tests}
Die Chaos-Tests sollen das Verhalten der Umgebung bei Ausfällen testen.
Dabei wurde von K6 Last erzeugt. Dies entspricht 50 Studierenden, die über einen Zeitraum von 10 Minuten Aufgaben in Codetask bearbeiten.
Währenddessen wurden verschiedene Ausfälle in den Umgebungen ausgelöst.
Für Kubernetes wurde dabei ein spezielles Tool namens Chaos Mesh verwendet, das kontrollierte Ausfälle im Cluster erzeugen kann.
Bei Docker wurden Skripte ausgeführt, die Fehler in der Umgebung auslösten, um Ausfälle zu nutzen. 
Ursprünglich war in diesen auch ein Test-Tool geplant, doch Docker erkannte, dass die Befehle für die Ausfälle über die Docker-API kamen, wodurch die im Docker Compose festgelegten Restart-Policys nicht verwendet wurden.

Die Chaostests wurden auf drei verschiedene Szenarien eingegrenzt.
Es gibt zwar deutlich mehr Arten von Chaostests, aber nicht alle eignen sich, um die Zuverlässigkeit von Docker Compose und Kubernetes fair zu vergleichen.
Ein Beispiel ist der Node-Failure-Test, der untersucht, was passiert, wenn ein Node ausfällt. Bei diesem Szenario ist das Ergebnis vorhersehbar, da bei Docker Compose ein Node-Ausfall einen kompletten Ausfall bedeutet.
„Partial Degradation”, bei dem der Großteil der Replicas eines Services ausfällt, ist wiederum mit Docker Compose nicht machbar, da dieser immer nur ein Replica pro Service hat.
Andere Chaostests sind nicht unbedingt Vergleiche zur Zuverlässigkeit, sondern eher Architektur-/Designfragen, wie z. B. Network Chaos, mit dem eher die Netzwerkqualität getestet wird.

Der erste Test ist der Instance-Failure: Dabei fällt ein Teil des Services aus. Bei Kubernetes entspricht dies einem Pod des Services, während bei Docker Compose bereits der gesamte Container des Services ausfällt.
Mit diesem Test soll die Fehlertoleranz der Umgebung bei einem Ausfall geprüft werden.

Beim zweiten Test, dem Service-Outage, fällt der komplette Service aus. Bei Kubernetes entspricht dies dem Ausfall aller Pods eines Service, während bei Docker Compose wieder der gesamte Container des Service ausfällt.
Mit diesem Test soll die Wiederherstellbarkeit der Umgebung getestet werden, d. h., wie lange es dauert, bis der Service wiederhergestellt ist und wieder verwendet werden kann.

Der dritte Test ist der Update-Test. Dieser testet, wie lange die Anwendung nicht erreichbar ist, wenn sich der Zustand durch ein Deployment ändert.
Mit diesem Test soll geprüft werden, wie verfügbar die Anwendungen bei einem Deployment bleiben. Bei Docker Compose wurde dazu der Webhook auf der VM ausgelöst, der das Deployment auslöst.
Bei Kubernetes wurde ein „kubectl rollout restart deployment” für den jeweiligen Service ausgelöst.

Alle Tests wurden jeweils nur mit der Haupt-API durchgeführt, da diese im Vergleich zu den anderen Services weniger Mehrwert bot. 
Die Ergebnisse wären ähnlich gewesen. Eine Ausnahme wäre der Update-Test, bei dem ein Unterschied entstanden wäre, wenn beispielsweise die AI-API aktualisiert wird.
Das Ergebnis ist dabei aber von vornherein vorhersehbar, da beim Docker Compose alle Services beim Update gleichzeitig heruntergefahren und wieder deployed werden, wodurch auch die Haupt-API ausfällt. Bei Kubernetes werden die Services dagegen einzeln deployed.